\section{Limitations and future work}

The study has several limitations. First, the sample is limited to three JSE Top 40 cross-sections and should not be generalised to the full South African equity market without further testing. The same firms recur, adjacent annual market windows overlap by six months and some financial statements are used at more than one date. The 123 instrument-date records therefore do not form 123 independent observations. The two Investec lines also represent one economic company group.

Second, the historical data are retrospectively reconstructed rather than fully certified as information available to an investor on each decision date. Shares and market capitalisation remain reconstructed, the December 2025 Investec allocation is interpolated, Reinet fundamentals are imputed and two Anglo American volume observations are imputed. JSE-only turnover does not capture global trading activity for internationally listed companies. Revenue definitions also differ across sectors. These issues can influence financial similarity even after the documented corrections.

Third, the UMAP settings were selected after examining June 2025, and eight clusters were retained from the earlier study rather than established through independent count selection. With 41 instruments and eight groups, several clusters contain few members, including two December singletons. Random-seed agreement assesses one source of variability but does not rule out sensitivity to the neighbourhood size, minimum distance, linkage, feature definition or preprocessing policy. The feature alternatives and cluster-count sensitivity show that the reported structure depends on these choices.

Fourth, matching groups across dates involves an interpretive choice. The Hungarian assignment maximises total continuing-member Jaccard similarity, but a one-to-one match cannot fully represent cluster splits and mergers. A company's label can change because its own profile changes, its peers change or the refitted embedding changes. Seed ARI does not provide a calibrated probability that an observed switch reflects economic change. The observed periods are also too few to establish a trend in stability.

Fifth, sector labels provide descriptive context but are not independently verified historical ground truth. The surrogate SHAP model explains cluster labels rather than the clustering algorithm itself, and its out-of-fold fidelity is limited or unavailable at some dates. The Mahalanobis reference assumes independent multivariate normal observations, which is not automatically appropriate for these financial cross-sections. Its nominal threshold is therefore a descriptive screen and does not establish statistically significant economic anomalies.

Finally, the study does not conduct portfolio backtesting, test subsequent return co-movement or establish that group changes predict risk or excess returns. Future work should verify point-in-time data availability and freeze inputs and settings before a new decision date. It should assess transition stability under input perturbations, alternative alignment objectives and preprocessing policies, and compare subsequent within-cluster return co-movement with sector-defined and matched random groups. Extending the analysis to more equities and longer periods would also allow the method to be assessed across different market conditions.
